\section{北方防线崩解时：城、路与人的最后选择}

\subsection{一二一一年以后，警报如何传递}

蒙古军进入金的北方边境后，中枢最关心的是哪一路敌军、哪座城和哪位将领；地方首先感到的却可能是驿站消息变慢、市场不敢远运、军队要求粮草、亲属从邻州涌来。警报沿道路传递，正如粮食和士兵沿道路移动。道路被截断时，朝廷即使仍发出命令，也难以确认谁已收到、谁有能力执行。

北方城池的得失不能只以“攻克”解释。城墙需要修补，守军需要轮换，城外的水源、农田与道路决定围城能维持多久。蒙古军的机动迫使金廷在许多地点同时分配兵力；金军与地方守备则利用城、河、山口和熟悉的路线抵抗。不同材料对战绩与责任的夸大都很明显，故本章不列没有互证的兵数。可以确定的是，战争逐渐破坏了使北方能够向中都输送人员、粮食和信息的条件。

\subsection{中都失去的是什么}

一二一五年中都失守，\eventref{JIN-1215-EVENT-104}{中都失守}。失去的不是一座象征性的首都而已。中都聚集官署、仓储、工匠、市场和通向燕云、河北、东北的路线；它被夺取后，留在原地的人、随朝廷南迁的人与进入城市的新力量，都要重新安排居住、供给和权利。城的陷落使旧有行政网络断开，却不使所有设施消失。后继政权仍须使用道路、市场和劳动力，只是使用者与受益者发生变化。

金廷转向河南，需要将官员、档案、军队与物资带入新的中枢。迁徙本身消耗粮食和车马，也会提高沿线地方的压力。新的行政中心依赖旧宋时期和金朝前期积累的城市与农业资源，却不能凭借这些资源自动获得地方信任。朝廷面对的是一连串时间不同的问题：今日要让官署开张，下一季要筹集军粮，几年后还要恢复能够支撑边防的税源。每个问题都可能同另一个问题争夺同一批人和物。

\subsection{围汴京的社会时间}

汴京受围时，军事与社会不再能分开。粮价、燃料、疫病、城门管理、消息和谣言都影响守城能力；而政治决定又会改变谁能出城、谁得到配给、谁被征入劳役。史书倾向于记录帝王、将领和忠臣，围城中的儿童、病者、临时工和无籍者则几乎没有自己的档案。承认这一点，不是把他们排除在历史之外，而是拒绝用虚构的细节替他们说话。

围城还使平日分散的财政问题暴露出来。仓中有多少粮、道路能否补给、军队是否得到饷给、城内市场能否运行，原本由多个机关和地区承担；敌军切断交通后，这些环节被压入同一地点。短期内可以靠征集和配给维持，时间越长，家庭储备、卫生和互助网络越容易耗尽。金末危机不是某项单独政策失败，而是城、路、仓和人同时被逼到极限。

\subsection{蔡州与短暂的合作}

哀宗退守蔡州，宋蒙共同攻城，最终金亡。这个结局在后来的朝代叙述中很容易被写成必然：金已衰亡，蒙古不可阻挡，南宋借机复仇。然而当时各方的目标并不相同。蒙古需要完成征服和控制北方，南宋希望取得河南方向的空间，金廷则试图在有限城池与外部交涉中延续。短期合作能把力量集中在蔡州，不能预先决定城破之后谁来管理粮区、城镇和道路。

蔡州陷落后，原先由金维持或限制的关系必须重新谈判。南宋北进的尝试和蒙古的行动很快显示，攻金并没有创造持久共同体。对于地方居民而言，旗帜再度变化可能先意味着新的征粮、军队与官吏；旧有的亲属、市场和寺院网络既可能提供庇护，也可能被新政权利用。金亡不是人的经历停止，而是他们要在更广阔的战争中重新选择。

\subsection{战争留下的文化后果}

战乱破坏书籍、碑刻和档案的保存条件，也改变谁有能力叙述过去。元修《金史》收集了大量金朝国史遗存，却以新政权的时间安排旧朝；宋方材料则以自身的战争、恢复或危机为中心。两者都让金末可见，又都使地方社会显得比宫廷和战场更沉默。城址、墓志、寺院与钱币材料能补正这种偏向，但保存是偶然的，发掘也不覆盖全部区域。

这层材料条件本身是晚金史的一部分。档案随朝廷南移或散失，家族在迁徙中丢失旧谱，寺院的营建和毁弃改变题记能否留下。今天看来像“史料空白”的地方，往往正是战争对行政和记忆造成的后果。面对它，历史写作的责任不是用连贯的末日图景填平，而是让读者看见哪些事可以确认、哪些人的经历还没有足够材料被说明。
